%% toc: P3 39 39.8
\chapter{Existential Pebble Games and Fixed-Subgraph Queries}
\label{sec:p03-w03-p03-061-existential}

Datalog programs and their extension by inequalities define queries that are
preserved under homomorphisms, respectively under one-to-one homomorphisms.
Preservation alone gives sufficient but not necessary conditions for
inexpressibility. This chapter supplies a tool that is both sound and
complete: existential finite-variable infinitary logics, which contain the
queries of bounded-variable Datalog, and an asymmetric pebble game in which
Spoiler plays only on the first structure. Duplicator wins this game exactly
when every existential \(k\)-variable sentence true in the first structure is
true in the second. The resulting criterion is then applied to the fixed
subgraph homeomorphism queries: for a star pattern the query is definable in
the existential logic with inequalities, a fact proved here by an explicit
Spoiler strategy combined with Menger's theorem.

\emph{Purpose.} To decide whether a query is expressible in the existential
infinitary logics, and hence whether it can be expressed by a Datalog or
Datalog\((\neq)\) program with a bounded number of variables.
\emph{Inputs.} A vocabulary, a class of finite structures, a Boolean query on
it, and for each \(k\) a pair of structures. \emph{Conclusions.} Either an
existential \(k\)-variable definition, or a proof of non-definability from
winning strategies of Duplicator in the existential pebble games.

The presentation follows Kolaitis's manuscript \emph{On the Expressive Power
of Logics on Finite Models}, dated 1 August 2003, Section~8, pages 69--79
\cite[Sect.~8, pp.~69--79]{Kolaitis2003ExpressivePowerNotes}. His chapter of
the same title in the 2007 volume \emph{Finite Model Theory and Its
Applications} is not the version followed; every locator refers to the 2003
manuscript. The manuscript states several results with a hint or a citation
only. All proofs below are written out in full in the annex, except where a
paragraph says explicitly that a clause is reported from the source without
proof. For the ordinary \(k\)-pebble game and the full finite-variable logics
see \cref{sec:p03-w03-p03-055-pebble}; for Ehrenfeucht--Fra\"iss\'e games see
\cref{sec:p03-ef-games}; for Datalog and conjunctive queries see
\cref{def:p09-datalog} and \cref{def:p09-conjunctive-query}.

\section{The existential infinitary logics}
\label{sec:p03-w03-p03-061-logics}

\begin{hbconvention}{def:p03-w03-p03-061-setting}{Vocabularies, partial homomorphisms and variables}
A vocabulary \(\sigma\) consists of relation symbols, each of positive arity,
and constant symbols \(c_1,\ldots,c_s\) (\(s\ge0\)). Structures
\(\mathbf A,\mathbf B\) have nonempty universes \(A,B\). The individual
variables are \(v_1,v_2,\ldots\); a \emph{term} is a variable or a constant
symbol. Put
\(C_{\mathbf A,\mathbf B}\defeq\{(c_1^{\mathbf A},c_1^{\mathbf B}),\ldots,
(c_s^{\mathbf A},c_s^{\mathbf B})\}\).

A \emph{partial homomorphism} from \(\mathbf A\) to \(\mathbf B\) is a
function \(f\) with \(\operatorname{dom}(f)\subseteq A\) and values in
\(B\) such that \(C_{\mathbf A,\mathbf B}\subseteq f\) and, for every
relation symbol \(R\) of arity \(r\) and all
\(a_1,\ldots,a_r\in\operatorname{dom}(f)\),
\(R^{\mathbf A}(a_1,\ldots,a_r)\) implies
\(R^{\mathbf B}(f(a_1),\ldots,f(a_r))\). It is \emph{one-to-one} when \(f\)
is injective. This is Kolaitis's requirement that \(a_i\mapsto b_i\) be a
homomorphism, respectively a one-to-one homomorphism, between the
substructures generated by the pebbled elements, which contain the constants
\cite[pp.~71--72]{Kolaitis2003ExpressivePowerNotes}. A homomorphism of
\(\sigma\)-structures is a partial homomorphism with domain \(A\).

For \(k\ge1\) put \(V_k\defeq\{v_1,\ldots,v_k\}\). A formula with at most
\(k\) distinct variables becomes, after a bijective renaming of its
variables, a formula whose variables lie in \(V_k\); the logics below are
taken in this normalized form. For \(\bar a=(a_1,\ldots,a_k)\in A^k\),
\(\mathbf A\models\varphi[\bar a]\) means satisfaction under
\(v_i\mapsto a_i\), and \(t^{\mathbf A}[\bar a]\) is the value of a term.
\end{hbconvention}

\begin{hbdefinition}{def:p03-w03-p03-061-logics}{The logics \(\exists\mathrm L^k_{\infty\omega}\) and \(\exists\mathrm L^k_{\infty\omega}(\neq)\)}
Let \(\sigma\) be a vocabulary and \(k\ge1\)
\cite[Def.~8.1, p.~69]{Kolaitis2003ExpressivePowerNotes}.
\begin{enumerate}
\item \(\exists\mathrm L^k_{\infty\omega}\) is the class of
      \(\mathrm L_{\infty\omega}\)-formulas with variables in \(V_k\) built
      from atomic formulas (equalities \(t_1=t_2\) included) by existential
      quantification \(\exists v_j\), conjunction \(\bigwedge\Psi\) and
      disjunction \(\bigvee\Psi\) of an arbitrary set \(\Psi\) of formulas.
      No negation and no universal quantifier occur. \(\exists\mathrm{FO}^k\)
      is the class of first-order (finitary) members of
      \(\exists\mathrm L^k_{\infty\omega}\).
\item \(\exists\mathrm L^k_{\infty\omega}(\neq)\) and
      \(\exists\mathrm{FO}^k(\neq)\) are defined in the same way, except
      that negated equalities \(t_1\neq t_2\) between terms are allowed as
      further basic formulas.
\item \(\exists\mathrm L^\omega_{\infty\omega}\defeq\bigcup_{k\ge1}
      \exists\mathrm L^k_{\infty\omega}\) and
      \(\exists\mathrm L^\omega_{\infty\omega}(\neq)\defeq\bigcup_{k\ge1}
      \exists\mathrm L^k_{\infty\omega}(\neq)\).
\end{enumerate}
The empty disjunction \(\bigvee\emptyset\) is written \(\bot\); it is false
in every structure and belongs to every \(\exists\mathrm L^k_{\infty\omega}\).
A query is \(\mathcal L\)-definable on a class \(\mathcal C\) of structures
when one \(\mathcal L\)-formula defines it on every member of \(\mathcal C\).
\end{hbdefinition}

\begin{hbexample}{def:p03-w03-p03-061-path3}{Paths with three variables}
Over \(\{E\}\) the \(\exists\mathrm{FO}^3\)-formula
\[
 (\exists v_3)\bigl(E(v_1,v_3)\wedge(\exists v_1)\bigl(v_1=v_3\wedge
 (\exists v_3)(E(v_1,v_3)\wedge E(v_3,v_2))\bigr)\bigr)
\]
holds of \((a_1,a_2)\) exactly when a walk of length \(3\) leads from
\(a_1\) to \(a_2\): the inner \(v_1\) is a copy of the first intermediate
node, after which \(v_3\) is free to name the second. Kolaitis writes this
formula with \(x,y,z\) and states that walks of every fixed length \(m\ge1\)
are \(\exists\mathrm{FO}^3\)-definable \cite[p.~70]{Kolaitis2003ExpressivePowerNotes};
\cref{thm:p03-w03-p03-061-walks} proves this for every \(m\). (The source
speaks of a ``path''; the formula does not prevent repeated nodes, so the
query it defines is the existence of a walk.)
\end{hbexample}

\begin{hbtheorem}{thm:p03-w03-p03-061-walks}{proved=annex:p03-w03-p03-061-walks}{Walks of every fixed length with three variables}
Over a vocabulary containing a binary relation symbol \(E\), put
\(\psi_1\defeq E(v_1,v_2)\) and
\[
 \psi_{m+1}\defeq(\exists v_3)\bigl(E(v_1,v_3)\wedge(\exists v_1)(v_1=v_3
 \wedge\psi_m)\bigr)\qquad(m\ge1).
\]
Each \(\psi_m\) is an \(\exists\mathrm{FO}^3\)-formula with free variables
among \(v_1,v_2\), and for every structure \(\mathbf A\) and
\(a_1,a_2,a_3\in A\), \(\mathbf A\models\psi_m[a_1,a_2,a_3]\) exactly when
there are \(e_0,\ldots,e_m\in A\) with \(e_0=a_1\), \(e_m=a_2\) and
\(E^{\mathbf A}(e_{r-1},e_r)\) for \(1\le r\le m\) (a walk of length \(m\)
from \(a_1\) to \(a_2\)).
\end{hbtheorem}

\begin{hbdefinition}{def:p03-w03-p03-061-datalog}{Datalog\((\neq)\) programs, stages and \(k\)-Datalog}
A \emph{Datalog\((\neq)\) program} \(\pi\) over a vocabulary \(\sigma\)
(the extensional symbols) is a finite set of rules
\[
 S(u_1,\ldots,u_n)\;\mathrel{{:}\!-}\;\beta_1,\ldots,\beta_q\qquad(q\ge0),
\]
where \(S\) is one of the intensional symbols \(S_1,\ldots,S_l\) (not in
\(\sigma\)), the \(u_i\) are terms, and each \(\beta_j\) is an atomic formula
\(R(\bar w)\) with \(R\in\sigma\cup\{S_1,\ldots,S_l\}\), an equality
\(w=w'\) or an inequality \(w\neq w'\) of terms. A \emph{Datalog} program has
no inequalities. The variables of the head need not occur in the body.

On a \(\sigma\)-structure \(\mathbf A\) the program determines an operator
\(\Phi\) on \(l\)-tuples \(\bar P=(P_1,\ldots,P_l)\) with
\(P_i\subseteq A^{n_i}\): \(\Phi(\bar P)_i\) is the set of tuples
\((u_1^{\mathbf A}[\beta],\ldots,u_{n_i}^{\mathbf A}[\beta])\), for rules
with head symbol \(S_i\) and assignments \(\beta\) of the rule's variables
into \(A\) under which every body literal holds in \((\mathbf A,\bar P)\). The
stages are \(\Phi^0\defeq(\emptyset,\ldots,\emptyset)\) and
\(\Phi^{m+1}\defeq\Phi(\Phi^m)\); \(\Phi^\infty\) is the least fixed point of
\(\Phi\). A program with a designated goal symbol \(S_i\) defines the query
\(\mathbf A\mapsto\Phi^\infty_i\). For \(k\ge1\), \(\pi\) is a
\(k\)-Datalog\((\neq)\) program when the body of every rule has at most \(k\)
distinct variables and its head has at most \(k\) distinct variables
\cite[p.~70]{Kolaitis2003ExpressivePowerNotes}; a rule as a whole may then
have more than \(k\) variables. This specializes \cref{def:p09-datalog}.
\end{hbdefinition}

\begin{hbexample}{def:p03-w03-p03-061-odd}{Non-2-colorability in 4-Datalog}
For a graph given by a symmetric relation \(E\), the program
\[
 O(x,y)\mathrel{{:}\!-} E(x,y);\qquad
 O(x,y)\mathrel{{:}\!-} E(x,z),E(z,w),O(w,y);\qquad
 Q\mathrel{{:}\!-} O(x,x)
\]
puts \((a,b)\) into \(O\) exactly when a walk of odd length leads from
\(a\) to \(b\) (induction on the stage for one inclusion, on the length of
the walk for the other), so the goal \(Q\) holds exactly when there is a
closed walk of odd length. A shortest closed walk of odd length is an odd
cycle, and a graph is 2-colorable exactly when it has no odd cycle, so \(Q\)
expresses non-2-colorability. The second rule has the four variables
\(x,y,z,w\); the program is in \(4\)-Datalog. Kolaitis prints the last body
atom of the second rule as \(O(z,y)\) \cite[p.~70]{Kolaitis2003ExpressivePowerNotes};
with that atom the rule adds walks of every length and the odd-cycle reading
fails, so the atom is taken here as \(O(w,y)\).
\end{hbexample}

\begin{hbtheorem}{thm:p03-w03-p03-061-datalog}{proved=annex:p03-w03-p03-061-datalog}{Bounded-variable Datalog is existential finite-variable}
Let \(\pi\) be a \(k\)-Datalog program over \(\sigma\) with intensional
symbols \(S_1,\ldots,S_l\) of arities \(n_1,\ldots,n_l\). For every \(i\) and
every map \(\mu\) from \(\{1,\ldots,n_i\}\) to the terms over \(V_k\):
\begin{enumerate}
\item for every \(m\ge0\) there is a formula \(\theta^m_{i,\mu}\), either
      \(\bot\) or a member of \(\exists\mathrm{FO}^k\), whose free variables
      occur among the terms \(\mu(1),\ldots,\mu(n_i)\), such that for every
      \(\sigma\)-structure \(\mathbf A\) and \(\bar a\in A^k\)
      \[
       \mathbf A\models\theta^m_{i,\mu}[\bar a]
       \iff
       \bigl(\mu(1)^{\mathbf A}[\bar a],\ldots,\mu(n_i)^{\mathbf A}[\bar a]\bigr)
       \in\Phi^m_i;
      \]
\item the formula \(\theta^\infty_{i,\mu}\defeq\bigvee_{m\ge0}\theta^m_{i,\mu}\)
      of \(\exists\mathrm L^k_{\infty\omega}\) satisfies the same equivalence
      with \(\Phi^\infty_i\) in place of \(\Phi^m_i\);
\item consequently every query defined by a \(k\)-Datalog program whose
      goal symbol has arity at most \(k\) is
      \(\exists\mathrm L^k_{\infty\omega}\)-definable on the class of all
      \(\sigma\)-structures: \(k\text{-Datalog}\subseteq
      \exists\mathrm L^k_{\infty\omega}\).
\end{enumerate}
The same holds for \(k\)-Datalog\((\neq)\) programs with
\(\exists\mathrm{FO}^k(\neq)\) and \(\exists\mathrm L^k_{\infty\omega}(\neq)\)
in place of \(\exists\mathrm{FO}^k\) and \(\exists\mathrm L^k_{\infty\omega}\).
\end{hbtheorem}

This is the Datalog form of Kolaitis's Theorems~8.2 and~8.3, for which he
gives a hint and refers to a later paper
\cite[pp.~70--71]{Kolaitis2003ExpressivePowerNotes}. The variable bound is
the point of the proof: a rule may mention up to \(2k\) variables, and the
translation recycles them through the substitution lemma
\cref{thm:p03-w03-p03-061-substitution}. His Theorem~8.2 itself concerns
systems of formulas, and is \cref{thm:p03-w03-p03-061-system} below.

\begin{hbdefinition}{def:p03-w03-p03-061-system}{Systems of positive existential formulas}
Let \(\sigma\) be a vocabulary, \(k\ge1\), and \(S_1,\ldots,S_l\) relation
symbols not in \(\sigma\), of arities \(n_1,\ldots,n_l\le k\). A
\emph{system of positive \(\exists\mathrm{FO}^k\)-formulas} is a list
\(\varphi_1,\ldots,\varphi_l\) of \(\exists\mathrm{FO}^k\)-formulas over
\(\sigma\cup\{S_1,\ldots,S_l\}\), the free variables of \(\varphi_i\) lying
among \(v_1,\ldots,v_{n_i}\) \cite[Thm~8.2, p.~70]{Kolaitis2003ExpressivePowerNotes}.
Its \emph{operator} on a \(\sigma\)-structure \(\mathbf A\) sends
\(\bar P=(P_1,\ldots,P_l)\), \(P_i\subseteq A^{n_i}\), to
\(\Phi(\bar P)\) with
\[
 \Phi(\bar P)_i\defeq\{(a_1,\ldots,a_{n_i}):(\mathbf A,\bar P)\models
 \varphi_i[a_1,\ldots,a_{n_i},\ldots]\},
\]
where the entries after \(a_{n_i}\) of the assignment are arbitrary. The
stages are \(\Phi^0\defeq(\emptyset,\ldots,\emptyset)\) and
\(\Phi^{m+1}\defeq\Phi(\Phi^m)\). Formulas of
\(\exists\mathrm{FO}^k\) contain no negation, so the \(S_j\) occur in them
only positively, which is what ``positive'' records. Systems of
\(\exists\mathrm{FO}^k(\neq)\)-formulas are defined in the same way.
\end{hbdefinition}

\begin{hbtheorem}{thm:p03-w03-p03-061-system}{proved=annex:p03-w03-p03-061-system}{Stages and least fixed points of positive existential systems}
For a system \(\varphi_1,\ldots,\varphi_l\) of positive
\(\exists\mathrm{FO}^k\)-formulas over \(\sigma\):
\begin{enumerate}
\item for every \(m\ge0\) and \(i\le l\) there is a formula
      \(\theta^m_i\), either \(\bot\) or a member of \(\exists\mathrm{FO}^k\),
      with free variables among \(v_1,\ldots,v_{n_i}\), that defines the
      component \(\Phi^m_i\) of the stage \(\Phi^m\) on the class of all
      \(\sigma\)-structures;
\item the operator \(\Phi\) is monotone, it has a least fixed point
      \(\Phi^\infty=\bigcup_m\Phi^m\) on every \(\sigma\)-structure, and
      \(\bigvee_{m\ge0}\theta^m_i\in\exists\mathrm L^k_{\infty\omega}\)
      defines \(\Phi^\infty_i\) on the class of all \(\sigma\)-structures.
\end{enumerate}
The same holds for systems of \(\exists\mathrm{FO}^k(\neq)\)-formulas, with
\(\exists\mathrm{FO}^k(\neq)\) and \(\exists\mathrm L^k_{\infty\omega}(\neq)\).
\end{hbtheorem}

This is Kolaitis's Theorem~8.2 \cite[Thm~8.2, p.~70]{Kolaitis2003ExpressivePowerNotes}.
Its second item prints \(\exists\mathrm{FO}^k\) for the components of the
least fixed point; the proof yields, and his concluding inclusion
\(k\text{-Datalog}\subseteq\exists\mathrm L^k_{\infty\omega}\) uses,
\(\exists\mathrm L^k_{\infty\omega}\). He obtains that inclusion by turning a
\(k\)-Datalog program into such a system and gives no details;
\cref{thm:p03-w03-p03-061-datalog} translates the rules of the program
directly instead.

\begin{hbtheorem}{thm:p03-w03-p03-061-proper}{proved=annex:p03-w03-p03-061-proper}{The inclusion of bounded-variable Datalog is proper}
Let \(\sigma\) contain a binary relation symbol \(E\), and let
\(\mathcal F\) be the class of finite \(\sigma\)-structures.
\begin{enumerate}
\item Every query defined on \(\mathcal F\) by a Datalog\((\neq)\) program is
      computable.
\item For every set \(S\) of prime numbers, the query \(Q_S\) ``some closed
      walk of \(E\) has a length in \(S\)'' is
      \(\exists\mathrm L^3_{\infty\omega}\)-definable on \(\mathcal F\); for
      some \(S\) it is not computable.
\end{enumerate}
Hence, on \(\mathcal F\), \(k\)-Datalog is properly contained in
\(\exists\mathrm L^\omega_{\infty\omega}\) and \(k\)-Datalog\((\neq)\) in
\(\exists\mathrm L^\omega_{\infty\omega}(\neq)\) for every \(k\ge1\), and for
\(k\ge3\) already in \(\exists\mathrm L^k_{\infty\omega}\), respectively
\(\exists\mathrm L^k_{\infty\omega}(\neq)\).
\end{hbtheorem}

Kolaitis states the proper inclusions with the reason that the infinitary
logic expresses non-computable queries
\cite[p.~71]{Kolaitis2003ExpressivePowerNotes}; the query \(Q_S\) is a
witness constructed here. A query is \emph{computable} when some algorithm,
given a finite \(\sigma\)-structure with universe \(\{1,\ldots,n\}\) listed
with its relations and constants, decides it.

\begin{hbtheorem}{thm:p03-w03-p03-061-substitution}{proved=annex:p03-w03-p03-061-substitution}{Substitution of terms for variables}
Let \(\mathcal L\) be one of \(\exists\mathrm L^k_{\infty\omega}\),
\(\exists\mathrm{FO}^k\), \(\exists\mathrm L^k_{\infty\omega}(\neq)\),
\(\exists\mathrm{FO}^k(\neq)\). For every \(\varphi\in\mathcal L\) and every
map \(\rho\) from \(\{1,\ldots,k\}\) to the terms over \(V_k\) there is
\(\varphi_\rho\in\mathcal L\) whose free variables occur among the terms
\(\rho(j)\) with \(v_j\) free in \(\varphi\), such that for every structure
\(\mathbf A\) and \(\bar a\in A^k\)
\[
 \mathbf A\models\varphi_\rho[\bar a]\iff
 \mathbf A\models\varphi\bigl[\rho(1)^{\mathbf A}[\bar a],\ldots,
 \rho(k)^{\mathbf A}[\bar a]\bigr].
\]
\end{hbtheorem}

For maps \(\rho\) into the variables this is the closure property that
Kolaitis singles out as the step to be established separately
\cite[p.~70]{Kolaitis2003ExpressivePowerNotes}.

\section{Existential pebble games}
\label{sec:p03-w03-p03-061-games}

The \(k\)-pebble game of \cref{def:p03-w03-p03-055-pebble} is symmetric:
Spoiler may play in either structure and the payoff asks for a partial
isomorphism. Moves of Spoiler in the second structure correspond to universal
quantifiers. Removing them, and weakening the payoff to a partial
homomorphism, gives a game matched to the existential logics
\cite[p.~71]{Kolaitis2003ExpressivePowerNotes}.

\begin{hbdefinition}{def:p03-w03-p03-061-game}{The \((\exists,k)\)- and \((\exists,\neq,k)\)-pebble games}
Let \(k\ge1\) and let \(\mathbf A,\mathbf B\) be \(\sigma\)-structures
\cite[Def.~8.4, pp.~71--72]{Kolaitis2003ExpressivePowerNotes}. Spoiler and
Duplicator each hold pebbles labelled \(1,\ldots,k\). In each move Spoiler
either places one of his pebbles that is not in use on an element of
\(A\), or removes one of his pebbles from \(A\); Duplicator answers by
placing her pebble with the same label on an element of \(B\), or by
removing her pebble with that label. A position determines the relation
\[
 f_p\defeq C_{\mathbf A,\mathbf B}\cup\{(a_j,b_j):\text{pebble }j
 \text{ lies on }a_j\in A\text{ and on }b_j\in B\}.
\]
Spoiler wins the \((\exists,k)\)-pebble game at a position when \(f_p\) is
not a partial homomorphism; he wins the \((\exists,\neq,k)\)-pebble game when
\(f_p\) is not a one-to-one partial homomorphism. Duplicator \emph{wins} a
game when she has a strategy under which Spoiler never wins, however long the
play.
\end{hbdefinition}

Kolaitis checks the payoff only once a pebble is down; checking it also at the
empty position changes no winner, since Spoiler can place a pebble at once.

\begin{hbdefinition}{def:p03-w03-p03-061-family}{Winning families}
A \emph{winning family} for the \((\exists,k)\)-pebble game (respectively the
\((\exists,\neq,k)\)-pebble game) on \(\mathbf A\) and \(\mathbf B\) is a
nonempty set \(\mathcal I\) of partial homomorphisms (respectively one-to-one
partial homomorphisms) from \(\mathbf A\) to \(\mathbf B\) such that
\cite[Def.~8.5, p.~72]{Kolaitis2003ExpressivePowerNotes}:
\begin{enumerate}
\item \(|f-C_{\mathbf A,\mathbf B}|\le k\) for every \(f\in\mathcal I\);
\item (closure under subfunctions) if \(g\in\mathcal I\) and
      \(C_{\mathbf A,\mathbf B}\subseteq f\subseteq g\), then
      \(f\in\mathcal I\);
\item (forth property up to \(k\)) if \(f\in\mathcal I\),
      \(|f-C_{\mathbf A,\mathbf B}|<k\) and \(a\in A\), then some
      \(g\in\mathcal I\) has \(f\subseteq g\) and
      \(a\in\operatorname{dom}(g)\).
\end{enumerate}
\end{hbdefinition}

\begin{hbtheorem}{thm:p03-w03-p03-061-strategy}{proved=annex:p03-w03-p03-061-strategy}{Strategies and families}
Duplicator wins the \((\exists,k)\)-pebble game (respectively the
\((\exists,\neq,k)\)-pebble game) on \(\mathbf A\) and \(\mathbf B\) if and
only if a winning family for that game exists. If \(j\le k\) and Duplicator
wins the game with \(k\) pebbles, she wins it with \(j\) pebbles.
\end{hbtheorem}

\begin{hbdefinition}{def:p03-w03-p03-061-preorder}{Existential \(k\)-variable preorders}
For \(\sigma\)-structures \(\mathbf A,\mathbf B\) and \(k\ge1\)
\cite[Def.~8.6, pp.~72--73]{Kolaitis2003ExpressivePowerNotes}:
\begin{enumerate}
\item \(\mathbf A\preceq^{\exists,k}_{\infty\omega}\mathbf B\) when every
      \(\exists\mathrm L^k_{\infty\omega}\)-sentence true in \(\mathbf A\) is
      true in \(\mathbf B\);
\item \(\mathbf A\preceq^{\exists,k}_{\omega\omega}\mathbf B\) when every
      \(\exists\mathrm{FO}^k\)-sentence true in \(\mathbf A\) is true in
      \(\mathbf B\);
\item for \(r\le k\), \(a_1,\ldots,a_r\in A\) and \(b_1,\ldots,b_r\in B\),
      \((\mathbf A,a_1,\ldots,a_r)\preceq^{\exists,k}_{\infty\omega}
      (\mathbf B,b_1,\ldots,b_r)\) when for every
      \(\exists\mathrm L^k_{\infty\omega}\)-formula \(\varphi\) with free
      variables among \(v_1,\ldots,v_r\),
      \(\mathbf A\models\varphi[a_1,\ldots,a_r]\) implies
      \(\mathbf B\models\varphi[b_1,\ldots,b_r]\).
\end{enumerate}
The relations \(\preceq^{\exists,\neq,k}_{\infty\omega}\) and
\(\preceq^{\exists,\neq,k}_{\omega\omega}\) are defined in the same way from
\(\exists\mathrm L^k_{\infty\omega}(\neq)\) and \(\exists\mathrm{FO}^k(\neq)\).
\end{hbdefinition}

\begin{hbtheorem}{thm:p03-w03-p03-061-game-logic}{proved=annex:p03-w03-p03-061-game-logic}{Existential pebble games characterize existential \(k\)-variable preservation}
Let \(k\ge1\) and let \(\mathbf A,\mathbf B\) be \(\sigma\)-structures. The
following are equivalent:
\begin{enumerate}
\item \(\mathbf A\preceq^{\exists,k}_{\infty\omega}\mathbf B\);
\item Duplicator wins the \((\exists,k)\)-pebble game on \(\mathbf A\) and
      \(\mathbf B\).
\end{enumerate}
If \(\mathbf B\) is finite, they are also equivalent to
\(\mathbf A\preceq^{\exists,k}_{\omega\omega}\mathbf B\). The same holds for
\(\preceq^{\exists,\neq,k}_{\infty\omega}\),
\(\preceq^{\exists,\neq,k}_{\omega\omega}\) and the
\((\exists,\neq,k)\)-pebble game.
\end{hbtheorem}

This is Kolaitis's Theorem~8.7, which he attributes to Kolaitis and Vardi and
states without proof \cite[Thm~8.7, p.~73]{Kolaitis2003ExpressivePowerNotes}.
The finiteness of \(\mathbf B\) enters only in the forth step of the proof,
where a conjunction indexed by the elements of \(B\) must be finitary.

\begin{hbtheorem}{thm:p03-w03-p03-061-homomorphism}{proved=annex:p03-w03-p03-061-homomorphism}{Preservation under homomorphisms}
If there is a homomorphism \(h\colon\mathbf A\to\mathbf B\), Duplicator wins
the \((\exists,k)\)-pebble game on \(\mathbf A\) and \(\mathbf B\) for every
\(k\), so every \(\exists\mathrm L^\omega_{\infty\omega}\)-sentence true in
\(\mathbf A\) is true in \(\mathbf B\). If \(h\) is one-to-one, the same holds
for the \((\exists,\neq,k)\)-pebble games and
\(\exists\mathrm L^\omega_{\infty\omega}(\neq)\).
\end{hbtheorem}

Kolaitis records these preservation properties as extensions of those of
Datalog and Datalog\((\neq)\) \cite[p.~71]{Kolaitis2003ExpressivePowerNotes}.

\begin{hbtheorem}{thm:p03-w03-p03-061-cliques}{proved=annex:p03-w03-p03-061-cliques}{The existential game is asymmetric}
Let \(\mathbf K_n\) be the loopless complete graph on \(n\) vertices and
\(k\ge1\). Spoiler wins the \((\exists,k+1)\)-pebble game on
\(\mathbf K_{k+1}\) and \(\mathbf K_k\), while Duplicator wins the
\((\exists,\neq,j)\)-pebble game, and hence the \((\exists,j)\)-pebble game,
on \(\mathbf K_k\) and \(\mathbf K_{k+1}\) for every \(j\ge1\).
\end{hbtheorem}

The first assertion is Kolaitis's example of the asymmetry
\cite[p.~72]{Kolaitis2003ExpressivePowerNotes}. By
\cref{thm:p03-w03-p03-061-game-logic}, an
\(\exists\mathrm{FO}^{k+1}\)-sentence (``there are \(k+1\) pairwise adjacent
vertices'') separates \(\mathbf K_{k+1}\) from \(\mathbf K_k\) in one
direction, while no existential sentence true in \(\mathbf K_k\) fails in
\(\mathbf K_{k+1}\).

\begin{hbtheorem}{thm:p03-w03-p03-061-criterion}{proved=annex:p03-w03-p03-061-criterion}{Game criterion for existential finite-variable definability}
Let \(\mathcal C\) be a class of finite \(\sigma\)-structures and \(Q\) a
Boolean query on \(\mathcal C\), invariant under isomorphism. The following
are equivalent:
\begin{enumerate}
\item \(Q\) is \(\exists\mathrm L^\omega_{\infty\omega}\)-definable on
      \(\mathcal C\);
\item there is \(k\ge1\) such that for all \(\mathbf A,\mathbf B\in\mathcal
      C\), if \(Q(\mathbf A)=1\) and Duplicator wins the
      \((\exists,k)\)-pebble game on \(\mathbf A\) and \(\mathbf B\), then
      \(Q(\mathbf B)=1\).
\end{enumerate}
For a fixed \(k\), item~2 holds exactly when \(Q\) is
\(\exists\mathrm L^k_{\infty\omega}\)-definable on \(\mathcal C\). The same
holds with \(\exists\mathrm L^\omega_{\infty\omega}(\neq)\),
\(\exists\mathrm L^k_{\infty\omega}(\neq)\) and the
\((\exists,\neq,k)\)-pebble games.
\end{hbtheorem}

This is Kolaitis's Corollary~8.8 \cite[Cor.~8.8, p.~73]{Kolaitis2003ExpressivePowerNotes},
with the fixed-\(k\) refinement made explicit. By
\cref{thm:p03-w03-p03-061-datalog}, non-definability in
\(\exists\mathrm L^\omega_{\infty\omega}(\neq)\) excludes every
Datalog\((\neq)\) program, whatever its number of variables.

\begin{hbmethod}{def:p03-w03-p03-061-method}{The method of existential pebble games}
Let \(\mathcal C\) be a class of finite \(\sigma\)-structures and \(Q\) a
Boolean query on \(\mathcal C\)
\cite[Method~8.9, p.~73]{Kolaitis2003ExpressivePowerNotes}.
\begin{enumerate}
\item (Soundness.) To show that \(Q\) is not
      \(\exists\mathrm L^\omega_{\infty\omega}\)-definable on \(\mathcal C\),
      produce for every \(k\ge1\) structures \(\mathbf A_k,\mathbf B_k\in
      \mathcal C\) with \(Q(\mathbf A_k)=1\), \(Q(\mathbf B_k)=0\), and a
      winning family for the \((\exists,k)\)-pebble game on \(\mathbf A_k\)
      and \(\mathbf B_k\).
\item (Completeness.) If \(Q\) is not
      \(\exists\mathrm L^\omega_{\infty\omega}\)-definable on \(\mathcal C\),
      such structures exist for every \(k\).
\item (Side conditions.) The pair must be produced for every \(k\); a pair
      for a single \(k\) shows only that \(Q\) is not
      \(\exists\mathrm L^k_{\infty\omega}\)-definable. Spoiler plays in
      \(\mathbf A_k\), the structure satisfying \(Q\). For
      \(\exists\mathrm L^\omega_{\infty\omega}(\neq)\) and Datalog\((\neq)\)
      the family must consist of one-to-one partial homomorphisms.
\end{enumerate}
Both items follow from \cref{thm:p03-w03-p03-061-criterion}. The source's
soundness clause names the game ``on \(\mathbf A\) and \(\mathbf B\)''; the
structures meant are \(\mathbf A_k\) and \(\mathbf B_k\).
\end{hbmethod}

\begin{hbexample}{def:p03-w03-p03-061-horn}{Horn formulas and the \((\exists,3)\)-pebble game}
Let \(\sigma\) consist of unary \(N_1,P_1\), binary \(N_2,P_2\) and ternary
\(N_3,P_3\), standing for Horn clauses \(\neg x\), \(x\), \(\neg x\vee\neg
y\), \(\neg x\vee y\), \(\neg x\vee\neg y\vee\neg z\) and \(\neg x\vee\neg
y\vee z\). Let \(\mathbf B\) have universe \(\{0,1\}\) and
\[
 N_1^{\mathbf B}=\{0\},\ P_1^{\mathbf B}=\{1\},\
 N_2^{\mathbf B}=\{0,1\}^2-\{(1,1)\},\
 P_2^{\mathbf B}=\{0,1\}^2-\{(1,0)\},
\]
\[
 N_3^{\mathbf B}=\{0,1\}^3-\{(1,1,1)\},\qquad
 P_3^{\mathbf B}=\{0,1\}^3-\{(1,1,0)\},
\]
the satisfying assignments of the six clause forms. A Horn formula
\(\varphi\) whose clauses have at most three literals is first normalized
as follows. If it contains an empty clause, replace it by \(x\wedge\neg x\)
for a fresh variable \(x\); otherwise leave its clauses unchanged. Add one
fresh unused variable to the variable set in either case, so the encoding
has a nonempty universe even for the empty conjunction. These operations
preserve satisfiability. The structure \(\mathbf A_\varphi\) has this
enlarged variable set as its universe, with one tuple for each normalized
clause (for instance
\((x,y,z)\in P_3^{\mathbf A_\varphi}\) when \(\neg x\vee\neg y\vee z\) is a
clause). This is Kolaitis's construction
\cite[proof of Prop.~8.11, p.~76]{Kolaitis2003ExpressivePowerNotes}; his
display gives the last relation the name \(N_3\) a second time, and it is the
relation \(P_3\).
\end{hbexample}

\begin{hbtheorem}{thm:p03-w03-p03-061-horn}{proved=annex:p03-w03-p03-061-horn}{Satisfiability as an existential pebble game}
With the notation of \cref{def:p03-w03-p03-061-horn}, the following are
equivalent for a Horn formula \(\varphi\) with at most three literals per
clause: (1) \(\varphi\) is satisfiable; (2) Duplicator wins the
\((\exists,k)\)-pebble game on \(\mathbf A_\varphi\) and \(\mathbf B\) for
every \(k\ge1\); (3) Duplicator wins the \((\exists,3)\)-pebble game on
\(\mathbf A_\varphi\) and \(\mathbf B\).
\end{hbtheorem}

The proof runs the marking algorithm for Horn satisfiability as the
3-Datalog program
\[
\begin{aligned}
 T(z)&\mathrel{{:}\!-}P_1(z);
 &P&\mathrel{{:}\!-}N_1(x),T(x);\\
 T(z)&\mathrel{{:}\!-}P_2(x,z),T(x);
 &P&\mathrel{{:}\!-}N_2(x,y),T(x),T(y);\\
 T(z)&\mathrel{{:}\!-}P_3(x,y,z),T(x),T(y);
 &P&\mathrel{{:}\!-}N_3(x,y,z),T(x),T(y),T(z).
\end{aligned}
\]
with goal \(P\), and transfers its goal from \(\mathbf A_\varphi\) to
\(\mathbf B\) by \cref{thm:p03-w03-p03-061-datalog} and
\cref{thm:p03-w03-p03-061-game-logic}.

The winner of the existential game can be computed by an induction that
collects the positions from which Spoiler forces a win. For \(\sigma\)-structures
\(\mathbf A,\mathbf B\), \(\bar a\in A^k\) and \(\bar b\in B^k\) put
\(f_{\bar a,\bar b}\defeq C_{\mathbf A,\mathbf B}\cup\{(a_1,b_1),\ldots,
(a_k,b_k)\}\), and write \(\bar a[j/a]\) for \(\bar a\) with its \(j\)-th entry
replaced by \(a\). The structure \(\mathbf A\oplus\mathbf B\) has as universe
the disjoint union of \(A\) and \(B\), unary relations \(D_1,D_2\) for the two
parts, and for each relation symbol \(R\) and constant symbol \(c\) of
\(\sigma\) the copies \(R^1,c^1\) interpreted as in \(\mathbf A\) and
\(R^2,c^2\) interpreted as in \(\mathbf B\).

\begin{hbtheorem}{thm:p03-w03-p03-061-spoiler-lfp}{proved=annex:p03-w03-p03-061-spoiler-lfp}{Determining the winner of the existential game}
Let \(k\ge1\) and let \(\mathbf A,\mathbf B\) be \(\sigma\)-structures. Let
\(\mathcal W_{\mathbf A,\mathbf B}\subseteq A^k\times B^k\) be the least fixed
point of the monotone operator \(\Gamma\) with \((\bar a,\bar b)\in\Gamma(T)\)
when \(f_{\bar a,\bar b}\) is not a partial homomorphism, or there are
\(j\le k\) and \(a\in A\) such that \((\bar a[j/a],\bar b[j/b])\in T\) for
every \(b\in B\) \cite[Thm~8.10, pp.~74--75]{Kolaitis2003ExpressivePowerNotes}.
\begin{enumerate}
\item Duplicator wins the \((\exists,k)\)-pebble game on \(\mathbf A\) and
      \(\mathbf B\) if and only if for every \(\bar a\in A^k\) there is
      \(\bar b\in B^k\) with \((\bar a,\bar b)\notin\mathcal W_{\mathbf
      A,\mathbf B}\).
\item There is a sentence of least fixed-point logic
      (\cref{def:p03-w03-p03-059-lfp}) that holds in \(\mathbf A\oplus\mathbf
      B\) exactly when Duplicator does not win the \((\exists,k)\)-pebble game
      on \(\mathbf A\) and \(\mathbf B\). For fixed \(\sigma\) and \(k\) this
      can be decided in time polynomial in \(|A|+|B|\) for finite
      \(\mathbf A,\mathbf B\).
\item For every finite \(\sigma\)-structure \(\mathbf B\) the \(k\)-Datalog
      program \(\pi_{\mathbf B}\) of \cref{def:p03-w03-p03-061-program-b}
      defines, on the class of all \(\sigma\)-structures \(\mathbf A\), the
      query ``Duplicator does not win the \((\exists,k)\)-pebble game on
      \(\mathbf A\) and \(\mathbf B\)''.
\end{enumerate}
The same holds for the \((\exists,\neq,k)\)-pebble game, with ``one-to-one
partial homomorphism'' in the definition of \(\Gamma\) and the
\(k\)-Datalog\((\neq)\) program obtained by adding the rules of
\cref{def:p03-w03-p03-061-program-b} for injectivity.
\end{hbtheorem}

\begin{hbdefinition}{def:p03-w03-p03-061-program-b}{The program \(\pi_{\mathbf B}\)}
Let \(\mathbf B\) be a finite \(\sigma\)-structure and \(k\ge1\). For each
\(\bar b\in B^k\) the program has a \(k\)-ary intensional symbol
\(T_{\bar b}\); \(\bar x\) is \(x_1,\ldots,x_k\), and a term is one of the
\(x_i\) or a constant symbol, read in \(\mathbf B\) at \(\bar b\) by
\(x_i\mapsto b_i\). The rules are
\begin{enumerate}
\item \(T_{\bar b}(\bar x)\mathrel{{:}\!-}t=t'\) for all terms \(t,t'\) with
      different values in \(\mathbf B\) at \(\bar b\);
\item \(T_{\bar b}(\bar x)\mathrel{{:}\!-}R(t_1,\ldots,t_r)\) for every
      relation symbol \(R\) of arity \(r\) and all terms with
      \(\mathbf B\not\models R(t_1,\ldots,t_r)\) at \(\bar b\);
\item for each \(j\le k\), the rule
      \(T_{\bar b}(\bar x)\mathrel{{:}\!-}T_{\bar b[j/c]}(x_1,\ldots,x_{j-1},y,
      x_{j+1},\ldots,x_k)\) for all \(c\in B\), the body being the list of
      these \(|B|\) atoms and \(y\) a new variable;
\item the goal rule \(\mathrm{Goal}\mathrel{{:}\!-}T_{\bar b}(\bar x)\) for
      all \(\bar b\in B^k\), the body being the list of these atoms.
\end{enumerate}
For the \((\exists,\neq,k)\)-pebble game add
\(T_{\bar b}(\bar x)\mathrel{{:}\!-}t\neq t'\) for all terms with the same
value in \(\mathbf B\) at \(\bar b\). This is Kolaitis's program, extended to
constant symbols \cite[pp.~74--75]{Kolaitis2003ExpressivePowerNotes}.
\end{hbdefinition}

Kolaitis gives only a sketch of \cref{thm:p03-w03-p03-061-spoiler-lfp}, and
for simplicity assumes that the vocabulary has no constant symbols. Every
body of \(\pi_{\mathbf B}\) has at most \(k\) variables (in the third rule the
head variable \(x_j\) is replaced by \(y\)) and every head at most \(k\), so
\(\pi_{\mathbf B}\) is a \(k\)-Datalog program, although its number of rules
grows with \(|B|^k\).

\begin{hbtheorem}{thm:p03-w03-p03-061-horn-complete}{proved=annex:p03-w03-p03-061-horn-complete}{A \(\mathsf P\)-complete existential game}
For the vocabulary and the structure \(\mathbf B\) of
\cref{def:p03-w03-p03-061-horn}, the problem ``given a finite
\(\sigma\)-structure \(\mathbf A\), does Duplicator win the
\((\exists,3)\)-pebble game on \(\mathbf A\) and \(\mathbf B\)?'' is in
\(\mathsf P\), and satisfiability of Horn formulas with at most three
literals per clause reduces to it in logarithmic space through
\(\varphi\mapsto\mathbf A_\varphi\). Since Horn 3-satisfiability is
\(\mathsf P\)-complete under logarithmic-space reductions, the problem is
\(\mathsf P\)-complete.
\end{hbtheorem}

This is Kolaitis's Proposition~8.11
\cite[Prop.~8.11, pp.~75--76]{Kolaitis2003ExpressivePowerNotes}. The
\(\mathsf P\)-completeness of Horn 3-satisfiability is taken from the
literature that he cites; the annex proves the membership and the reduction.
Kolaitis also reports from Kolaitis and Panttaja that the problem is
\(\mathsf{EXPTIME}\)-complete when \(k\) is part of the input
\cite[Thm~8.12, p.~77]{Kolaitis2003ExpressivePowerNotes}; that result is
reported here without proof.

\section{Fixed subgraph homeomorphism queries}
\label{sec:p03-w03-p03-061-fish}

\begin{hbconvention}{def:p03-w03-p03-061-pattern}{Pattern graphs and disjoint paths}
Directed graphs are pairs \((V,E)\) with \(E\subseteq V\times V\). A
pattern graph \(\mathbf H=(V_H,E_H)\) is finite, has at least one node,
has no loops, and has no isolated nodes, so it has at least one edge; Kolaitis and Vardi adopt the last restriction explicitly
\cite[p.~22]{KolaitisVardi1995DatalogManuscript}, and
\cref{def:p03-w03-p03-061-edgeless} shows that the dichotomy below needs it.
Simple paths \(p\) and \(p'\) are \emph{node-disjoint} when every node they
share is an endpoint of both, following the reading
\cite[p.~22, n.~3]{KolaitisVardi1995DatalogManuscript}. Kolaitis's
``pairwise node-disjoint simple paths'' is taken in this sense.
\end{hbconvention}

\begin{hbdefinition}{def:p03-w03-p03-061-fish}{Homeomorphisms and the queries \(\mathrm{FiSH}(\mathbf H)\)}
Let \(\mathbf H\) and \(\mathbf G\) be directed graphs
\cite[Defs~8.13--8.14, p.~77]{Kolaitis2003ExpressivePowerNotes}. A
\emph{homeomorphism} \(h\colon\mathbf H\rightsquigarrow\mathbf G\) is a
one-to-one map from the nodes of \(\mathbf H\) to the nodes of \(\mathbf G\)
together with, for each edge \(e=(u,w)\) of \(\mathbf H\), a simple path
\(p_e\) of \(\mathbf G\) from \(h(u)\) to \(h(w)\), the paths being pairwise
node-disjoint. The \emph{fixed subgraph homeomorphism query with pattern
\(\mathbf H\)}, \(\mathrm{FiSH}(\mathbf H)\), asks of a directed graph
\(\mathbf G\) and a one-to-one map from the nodes of \(\mathbf H\) to the
nodes of \(\mathbf G\) whether some homeomorphism extends the map.

As a query on structures: let \(\sigma_{\mathbf H}\) consist of a binary
\(E\) and a constant \(d_v\) for each node \(v\) of \(\mathbf H\), and let
\(\mathcal F_{\mathbf H}\) be the class of finite
\(\sigma_{\mathbf H}\)-structures in which the constants denote pairwise
distinct elements. \(\mathrm{FiSH}(\mathbf H)(\mathbf G)=1\) when the map
\(v\mapsto d_v^{\mathbf G}\) extends to a homeomorphism.
\end{hbdefinition}

\begin{hbexample}{def:p03-w03-p03-061-fish-examples}{Disjoint paths and cycles through given nodes}
\cite[Exs~8.15--8.16, pp.~77--78]{Kolaitis2003ExpressivePowerNotes}
\begin{enumerate}
\item If \(\mathbf H\) has nodes \(s_1,s_2,t_1,t_2\) and edges
      \((s_1,t_1),(s_2,t_2)\), \(\mathrm{FiSH}(\mathbf H)\) is the
      \emph{2-disjoint paths} query: given \(\mathbf G\) and four nodes
      \(s_1',s_2',t_1',t_2'\), are there node-disjoint simple paths from
      \(s_1'\) to \(t_1'\) and from \(s_2'\) to \(t_2'\)? With \(m\)
      pairwise vertex-disjoint directed edges, with all \(2m\) endpoints
      distinct, it is the \(m\)-disjoint paths query.
\item If \(\mathbf H=\mathbf C_m\) is a directed cycle with \(m\ge3\) nodes,
      \(\mathrm{FiSH}(\mathbf C_m)\) asks whether \(\mathbf G\) has a simple
      cycle passing through the \(m\) given nodes (in the cyclic order fixed
      by the map).
\end{enumerate}
\end{hbexample}

\begin{hbdefinition}{def:p03-w03-p03-061-star}{Star graphs}
A \emph{star graph} is a directed graph consisting of one \emph{source} node
and edges from it to the other nodes, or of one \emph{sink} node and edges
from the other nodes to it \cite[Def.~8.17, p.~78]{Kolaitis2003ExpressivePowerNotes}.
Under \cref{def:p03-w03-p03-061-pattern} a star pattern has a centre \(r\),
leaves \(w_1,\ldots,w_m\) with \(m\ge1\), and edges \((r,w_i)\) for all
\(i\) (an out-star) or \((w_i,r)\) for all \(i\) (an in-star).
\end{hbdefinition}

Fortune, Hopcroft and Wyllie classified the computational complexity of all
these queries: \(\mathrm{FiSH}(\mathbf H)\) is in \(\mathsf P\) when
\(\mathbf H\) is a star graph and \(\mathsf{NP}\)-complete otherwise, as
Kolaitis reports \cite[Thm~8.18, p.~78]{Kolaitis2003ExpressivePowerNotes}. By
Ladner's theorem, \(\mathsf{NP}\) contains problems of intermediate
complexity if \(\mathsf P\neq\mathsf{NP}\); the classification says that no
\(\mathrm{FiSH}(\mathbf H)\) is one of them, and it is a proper dichotomy only
under that hypothesis. Kolaitis then states a dichotomy for descriptive
complexity that needs no such hypothesis \cite[Thm~8.19, p.~79]{Kolaitis2003ExpressivePowerNotes}:
if \(\mathbf H\) is a star graph, \(\mathrm{FiSH}(\mathbf H)\) is
\(\exists\mathrm L^\omega_{\infty\omega}(\neq)\)-definable, and in fact
definable in Datalog\((\neq)\); if \(\mathbf H\) is not a star graph, it is
not \(\exists\mathrm L^\omega_{\infty\omega}(\neq)\)-definable. The first
clause, with its Datalog\((\neq)\) refinement, is proved below
(\cref{thm:p03-w03-p03-061-star,thm:p03-w03-p03-061-star-datalog}). The
second clause is proved in two steps. Every non-star pattern contains one of
three basic patterns, and non-definability passes from a pattern to every
larger one (\cref{thm:p03-w03-p03-061-basic,thm:p03-w03-p03-061-transfer});
the three basic queries are then shown not to be
\(\exists\mathrm L^\omega_{\infty\omega}(\neq)\)-definable in
\cref{sec:p03-w03-p03-061-switches}, and
\cref{thm:p03-w03-p03-061-dichotomy} assembles the parts. Kolaitis's own
argument for the second clause is a hint: for every \(k\) there are graphs
\(\mathbf A_k\) (two long disjoint paths) satisfying the 2-disjoint paths
query and \(\mathbf B_k\) failing it, on which Duplicator wins the
\((\exists,\neq,k)\)-pebble game, and \(\mathbf B_k\) is extracted from the
reduction of Fortune, Hopcroft and Wyllie that shows the 2-disjoint paths
query \(\mathsf{NP}\)-hard \cite[p.~79]{Kolaitis2003ExpressivePowerNotes}.
The construction carried out below is that of Kolaitis and Vardi
\cite[Thms~6.6--6.7, pp.~34--39]{KolaitisVardi1995DatalogManuscript}. Their
manuscript states the lower bound for a logic it writes \(L^\omega\); the
property it verifies for Duplicator's answers is that the map from pebbled
nodes to pebbled nodes is a one-to-one homomorphism
\cite[p.~38]{KolaitisVardi1995DatalogManuscript}, which is the payoff of the
\((\exists,\neq,k)\)-pebble game of \cref{def:p03-w03-p03-061-game}.

\begin{hbtheorem}{thm:p03-w03-p03-061-menger}{proved=annex:p03-w03-p03-061-menger}{Menger's theorem for directed graphs, fan form}
Let \(\mathbf D=(V,E)\) be a finite directed graph.
\begin{enumerate}
\item For \(X,Y\subseteq V\), the largest number of pairwise vertex-disjoint
      paths from \(X\) to \(Y\) equals the least size of a set
      \(Z\subseteq V\) meeting every path from \(X\) to \(Y\).
\item Let \(s\in V\) and \(T\subseteq V-\{s\}\) with \(|T|=m\). Either there
      are \(m\) paths from \(s\) to \(T\), each meeting \(T\) only in its
      last node and any two sharing only \(s\), or there is a set
      \(Z\subseteq V-\{s\}\) with \(|Z|\le m-1\) meeting every walk from
      \(s\) to a node of \(T\).
\end{enumerate}
\end{hbtheorem}

\begin{hbtheorem}{thm:p03-w03-p03-061-star}{proved=annex:p03-w03-p03-061-star}{Star patterns are existentially definable with inequalities}
Let \(\mathbf H\) be a star pattern with \(m\) edges, and let
\(\mathbf G,\mathbf G'\in\mathcal F_{\mathbf H}\). If
\(\mathrm{FiSH}(\mathbf H)(\mathbf G)=1\) and
\(\mathrm{FiSH}(\mathbf H)(\mathbf G')=0\), Spoiler wins the
\((\exists,\neq,k)\)-pebble game on \(\mathbf G\) and \(\mathbf G'\) for
every \(k\ge m+1\). Consequently \(\mathrm{FiSH}(\mathbf H)\) is
\(\exists\mathrm L^{m+1}_{\infty\omega}(\neq)\)-definable on
\(\mathcal F_{\mathbf H}\).
\end{hbtheorem}

Spoiler walks along the \(m\) disjoint paths of \(\mathbf G\) one at a time,
with two pebbles leapfrogging. Duplicator's answers trace a walk in
\(\mathbf G'\) that must cross a small separating set given by Menger's
theorem; Spoiler leaves a pebble at the crossing, and after \(m\) paths two
of his \(m\) distinct anchors have the same image.

The Datalog\((\neq)\) refinement follows Kolaitis and Vardi's proof of their
Theorem~6.1 \cite[Thm~6.1, pp.~23--25]{KolaitisVardi1995DatalogManuscript}.
It strengthens the query so that it can be computed by induction on the
number of paths: besides the targets, a list of nodes to be avoided is
carried along, and the avoided list grows by one at each level of the
recursion. Kolaitis and Vardi write \(Q_{k,l}\); the first index is \(j\)
here, since \(k\) counts pebbles in this chapter.

\begin{hbdefinition}{def:p03-w03-p03-061-avoiding}{Avoiding fans and the program \(\pi_m\)}
Let \(\mathbf G=(V,E)\) be a directed graph, \(j\ge1\) and \(l\ge0\). For
nodes \(s,s_1,\ldots,s_j,t_1,\ldots,t_l\) of \(\mathbf G\), the relation
\(Q_{j,l}(s,s_1,\ldots,s_j,t_1,\ldots,t_l)\) holds when there are simple
paths \(p_1,\ldots,p_j\), \(p_i\) from \(s\) to \(s_i\), any two of which
share only \(s\), and none of which contains any of the nodes
\(t_1,\ldots,t_l\) \cite[pp.~24--25]{KolaitisVardi1995DatalogManuscript}.
A path \emph{contains} its endpoints, so \(Q_{j,l}\) fails when \(s\) or
some \(s_i\) is among the \(t\)'s.

For \(m\ge1\), the Datalog\((\neq)\) program \(\pi_m\) over \(\{E\}\) has the
intensional symbols \(R_j\) of arity \(m+1\) for \(1\le j\le m\) (standing
for \(Q_{j,m-j}\)). With \(l\defeq m-j\), \(\bar y\defeq y_1,\ldots,y_{j-1}\),
\(\bar t\defeq t_1,\ldots,t_l\) and \(\bar t\) written \(t_1,\ldots,t_l\) in
the atoms, its rules are, for each \(j\),
\[
\begin{aligned}
 R_j(s,\bar y,x,\bar t)&\mathrel{{:}\!-}E(s,x),\ s\neq t_1,\ldots,s\neq t_l,\\
 &\qquad x\neq t_1,\ldots,x\neq t_l,\ \Lambda_j(s,\bar y,x,\bar t);\\
 R_j(s,\bar y,x,\bar t)&\mathrel{{:}\!-}R_j(s,\bar y,w,\bar t),\ E(w,x),\\
 &\qquad x\neq t_1,\ldots,x\neq t_l,\ \Lambda_j(s,\bar y,x,\bar t),
\end{aligned}
\]
where \(\Lambda_1\) is empty and, for \(j\ge2\),
\(\Lambda_j(s,\bar y,x,\bar t)\) is the atom
\[
 R_{j-1}(s,y_1,\ldots,y_{j-1},x,t_1,\ldots,t_l):
\]
the target \(y_{j-1}\)
of the inner fan is its last target, and \(x\) joins the avoided list. For a
star pattern \(\mathbf H\) with centre \(r\) and leaves \(w_1,\ldots,w_m\),
the program \(\pi_{\mathbf H}\) is \(\pi_m\) together with the goal rule
\(\mathrm{Goal}\mathrel{{:}\!-}R_m(d_r,d_{w_1},\ldots,d_{w_m})\) if
\(\mathbf H\) is an out-star; for an in-star, every atom \(E(u,u')\) of
\(\pi_m\) is replaced by \(E(u',u)\).
\end{hbdefinition}

\begin{hbtheorem}{thm:p03-w03-p03-061-star-datalog}{proved=annex:p03-w03-p03-061-star-datalog}{Star patterns are Datalog\((\neq)\)-definable}
Let \(m\ge1\), let \(\mathbf G\) be a finite directed graph and let
\(R_j^\infty\) be the least fixed point components of \(\pi_m\) on
\(\mathbf G\). For \(1\le j\le m\), \(l=m-j\), and all nodes
\(s,s_1,\ldots,s_j,t_1,\ldots,t_l\) with \(s,s_1,\ldots,s_j\) pairwise
distinct,
\[
 (s,s_1,\ldots,s_j,t_1,\ldots,t_l)\in R_j^\infty
 \iff Q_{j,l}(s,s_1,\ldots,s_j,t_1,\ldots,t_l).
\]
Consequently, for every star pattern \(\mathbf H\) with \(m\) edges, the
goal of \(\pi_{\mathbf H}\) defines \(\mathrm{FiSH}(\mathbf H)\) on
\(\mathcal F_{\mathbf H}\); \(\pi_{\mathbf H}\) is an
\((m+2)\)-Datalog\((\neq)\) program.
\end{hbtheorem}

\begin{hbwarning}{def:p03-w03-p03-061-kv-program}{The avoided nodes must be excluded from the walk}
Kolaitis and Vardi compute \(Q_{k,l}\) through an auxiliary query
\(Q'_{k,l}\) asking for a walk from \(s\) to \(s_k\) every later node \(w\)
of which satisfies
\[
 Q_{k-1,l+1}(s,s_1,\ldots,s_{k-1},w,t_1,\ldots,t_l);
\]
neither this description nor the two printed rules for \(Q'_{k,l}\) require
the walk to avoid \(t_1,\ldots,t_l\)
\cite[p.~25]{KolaitisVardi1995DatalogManuscript}, while their Menger step
uses a walk that avoids them. Without that requirement the claimed
equivalence fails: in the graph with edges \((s,s_1),(s,t),(t,s_2)\) the
printed rules derive \(Q'_{2,1}(s,s_1,s_2,t)\), through the walk
\(s,t,s_2\), although the only path from \(s\) to \(s_2\) passes through
\(t\), so \(Q_{2,1}(s,s_1,s_2,t)\) fails. The inequalities
\(x\neq t_1,\ldots,x\neq t_l\) in the rules of \(\pi_m\) supply the missing
requirement.
\end{hbwarning}

\begin{hbtheorem}{thm:p03-w03-p03-061-fish-preserved}{proved=annex:p03-w03-p03-061-fish-preserved}{Homeomorphism queries are preserved by one-to-one homomorphisms}
For every pattern \(\mathbf H\), if \(\mathrm{FiSH}(\mathbf H)(\mathbf G)=1\)
and \(g\colon\mathbf G\to\mathbf G'\) is a one-to-one homomorphism of
\(\sigma_{\mathbf H}\)-structures, then
\(\mathrm{FiSH}(\mathbf H)(\mathbf G')=1\).
\end{hbtheorem}

\begin{hbwarning}{def:p03-w03-p03-061-preservation-limit}{Preservation does not decide these queries}
Every \(\exists\mathrm L^\omega_{\infty\omega}(\neq)\)-definable query is
preserved by one-to-one homomorphisms
(\cref{thm:p03-w03-p03-061-homomorphism}), and by
\cref{thm:p03-w03-p03-061-fish-preserved} so is every
\(\mathrm{FiSH}(\mathbf H)\), star or not. Preservation therefore cannot
separate the two sides of the dichotomy; a game argument is needed for the
non-star side, as Kolaitis notes \cite[p.~79]{Kolaitis2003ExpressivePowerNotes}.
\end{hbwarning}

\begin{hbwarning}{def:p03-w03-p03-061-edgeless}{Isolated nodes and the literal dichotomy}
If isolated nodes were allowed, a pattern with two nodes and no edge would
not be a star graph, yet \(\mathrm{FiSH}\) of it holds on every member of
\(\mathcal F_{\mathbf H}\) (the empty family of paths extends the given
one-to-one map) and is defined there by the atomic sentence \(d_v=d_v\) of
\(\exists\mathrm L^1_{\infty\omega}\). The second clause of the dichotomy is
therefore false for such patterns, and the restriction of
\cref{def:p03-w03-p03-061-pattern} is needed. The same holds for the pattern with
no nodes, whose query holds everywhere and is defined by the empty
conjunction.
\end{hbwarning}

\begin{hbtheorem}{thm:p03-w03-p03-061-single-loop}{proved=annex:p03-w03-p03-061-single-loop}{The single-loop boundary}
Consider the extension in which the single edge of a one-node loop pattern
must be realized by a directed simple cycle of positive length through its
named image \(c\). A self-loop in the input graph counts as a cycle of
length one. On finite directed graphs this query is defined by the
two-variable Datalog program
\[
\begin{aligned}
 R(x)&\mathrel{{:}\!-}E(c,x),\\
 R(y)&\mathrel{{:}\!-}R(x),E(x,y),
\end{aligned}
\]
with Boolean output \(R(c)\). It is consequently definable in
\(\exists\mathrm L^2_{\infty\omega}\). The one-node loop pattern is not
an in-star or out-star with distinct centre and leaves. Thus the non-star
clause of \cref{thm:p03-w03-p03-061-dichotomy} cannot simply be extended
to patterns with loops under this cycle convention.
\end{hbtheorem}

This is a direct boundary example, not an extension of the source's
classification. It also explains why replacing a loop by a path of length
zero would change the query: that alternative convention makes the
single-loop query true on every input. Neither convention is silently
included in \cref{def:p03-w03-p03-061-pattern}.

The non-star side reduces to three small patterns, which Kolaitis and Vardi
list \cite[pp.~27--28]{KolaitisVardi1995DatalogManuscript}:
\(\mathbf H_1\) with distinct nodes \(s_1,s_2,s_3,s_4\) and edges
\((s_1,s_2),(s_3,s_4)\), whose query is the 2-disjoint paths query;
\(\mathbf H_2\) with distinct nodes \(s_1,s_2,s_3\) and edges
\((s_1,s_2),(s_2,s_3)\), asking for a simple path from \(s_1\) to \(s_3\)
through \(s_2\); and \(\mathbf H_3\) with distinct nodes \(s_1,s_2\) and
edges \((s_1,s_2),(s_2,s_1)\), asking for a simple cycle through \(s_1\) and
\(s_2\).

\begin{hblemma}{thm:p03-w03-p03-061-basic}{proved=annex:p03-w03-p03-061-basic}{Every non-star pattern contains a basic pattern}
A pattern graph that is not a star graph has a subgraph isomorphic to
\(\mathbf H_1\), \(\mathbf H_2\) or \(\mathbf H_3\).
\end{hblemma}

\begin{hbtheorem}{thm:p03-w03-p03-061-transfer}{proved=annex:p03-w03-p03-061-transfer}{Non-definability passes to larger patterns}
Let \(\mathbf F_1\) and \(\mathbf F_2\) be pattern graphs and \(\mathbf F_1\)
isomorphic to a subgraph of \(\mathbf F_2\), and let \(k\ge1\). If
\(\mathrm{FiSH}(\mathbf F_1)\) is not
\(\exists\mathrm L^k_{\infty\omega}(\neq)\)-definable on
\(\mathcal F_{\mathbf F_1}\), then \(\mathrm{FiSH}(\mathbf F_2)\) is not
\(\exists\mathrm L^k_{\infty\omega}(\neq)\)-definable on
\(\mathcal F_{\mathbf F_2}\). Consequently, if
\(\mathrm{FiSH}(\mathbf H_1)\), \(\mathrm{FiSH}(\mathbf H_2)\) and
\(\mathrm{FiSH}(\mathbf H_3)\) are not
\(\exists\mathrm L^\omega_{\infty\omega}(\neq)\)-definable, then no
\(\mathrm{FiSH}(\mathbf H)\) with a non-star pattern \(\mathbf H\) is.
\end{hbtheorem}

This is Kolaitis and Vardi's Lemma~6.3
\cite[Lemma~6.3, pp.~28--29]{KolaitisVardi1995DatalogManuscript}. Their
construction adds to the graphs a copy of the edges of \(\mathbf F_2\) that
are not in \(\mathbf F_1\), and for the step showing that the enlarged
negative instance still fails the query they refer to an induction in
Fortune, Hopcroft and Wyllie. The annex uses the same construction, the
added edges joining distinguished nodes, and proves that step directly:
a path for an edge of \(\mathbf F_1\) cannot pass through a distinguished
node other than its own endpoints.

\section{The non-star side: switches and existential games}
\label{sec:p03-w03-p03-061-switches}

This section proves that \(\mathrm{FiSH}(\mathbf H_1)\),
\(\mathrm{FiSH}(\mathbf H_2)\) and \(\mathrm{FiSH}(\mathbf H_3)\) are not
\(\exists\mathrm L^\omega_{\infty\omega}(\neq)\)-definable, by the method of
\cref{def:p03-w03-p03-061-method}. For every \(k\) the negative instance
\(\mathbf B_k\) is the graph that the reduction of Fortune, Hopcroft and
Wyllie assigns to the unsatisfiable formula \(\varphi_k\) whose clauses are
all \(2^k\) choices of one literal for each of \(k\) variables; the positive
instance \(\mathbf A_k\) consists of two disjoint paths whose lengths are those
of the two paths that a satisfying assignment would produce in
\(\mathbf B_k\). Duplicator answers a pebble on \(\mathbf A_k\) with the
vertex at the same place on one of these would-be paths, choosing the truth
values that decide the place only when a pebble forces her to; with \(k\)
pebbles she never meets all \(k\) literals of a clause already decided.

The plan and the graphs are those of Kolaitis and Vardi
\cite[pp.~29--39]{KolaitisVardi1995DatalogManuscript}. Their manuscript
describes the switch through six named paths and refers to figures (its
Figures~1--6) that are not part of the copy read. The switch below is read
from the reproduction of Fortune, Hopcroft and Wyllie's figure in
Bang-Jensen and Gutin's \emph{Digraphs}, in the corrected version of its
first edition dated 15 August 2007 that the authors make freely available
\cite[Fig.~9.1, p.~479]{BangJensenGutin2007DigraphsCorrected}, and the way
the columns, clause paths and switches are joined is fixed here as stated in
\cref{def:p03-w03-p03-061-fhw-graph}, following the four steps of the
manuscript's description \cite[p.~31]{KolaitisVardi1995DatalogManuscript}.
Kolaitis and Vardi describe Duplicator's moves through an auxiliary pebble
game on formulas \cite[Def.~6.5, p.~33]{KolaitisVardi1995DatalogManuscript};
here her strategy is given as a winning family
(\cref{thm:p03-w03-p03-061-family}) that records the same choices.

\begin{hbdefinition}{def:p03-w03-p03-061-switch}{The switch}
The \emph{switch} \(\mathbf S\) is the directed graph with the twenty
vertices
\[
 a,b,c,d,e,f,g,h,\quad \alpha_1,\alpha_2,\alpha_3,\alpha_4,\quad
 \alpha'_1,\alpha'_2,\alpha'_3,\alpha'_4,\quad
 \beta_1,\beta_2,\beta'_1,\beta'_2
\]
and the twenty-four edges
\[
\begin{array}{llllll}
 c\to\alpha_1, & \alpha_1\to\alpha_2, & \alpha_2\to\alpha_3, &
 \alpha_3\to\alpha_4, & \alpha_4\to a, & \alpha_2\to h,\\
 b\to\alpha_3, & \alpha_4\to\beta'_1, & g\to\beta_1, &
 \beta_1\to\beta_2, & \beta_2\to\alpha_1, & \beta_2\to d,\\
 c\to\alpha'_1, & \alpha'_1\to\alpha'_2, & \alpha'_2\to\alpha'_3, &
 \alpha'_3\to\alpha'_4, & \alpha'_4\to a, & \alpha'_2\to f,\\
 b\to\alpha'_3, & \alpha'_4\to\beta_1, & e\to\beta'_1, &
 \beta'_1\to\beta'_2, & \beta'_2\to\alpha'_1, & \beta'_2\to d.
\end{array}
\]
The \emph{mirror} \(\mu\) is the permutation of the vertices that exchanges
\(\alpha_i\leftrightarrow\alpha'_i\), \(\beta_i\leftrightarrow\beta'_i\),
\(g\leftrightarrow e\), \(h\leftrightarrow f\) and fixes \(a,b,c,d\); it maps
the first two rows of edges onto the last two, so it is an automorphism of
\(\mathbf S\). The vertices \(b,c,e,g\) have no incoming edge (\emph{entry}
vertices) and \(a,d,f,h\) have no outgoing edge (\emph{exit} vertices). A
path \emph{passes through} \(\mathbf S\) when it is a simple path of
\(\mathbf S\) from an entry vertex to an exit vertex
\cite[p.~30]{KolaitisVardi1995DatalogManuscript}. Six such paths are
singled out:
\[
\begin{array}{ll}
 p(c,a)=c\,\alpha_1\,\alpha_2\,\alpha_3\,\alpha_4\,a, &
 q(c,a)=c\,\alpha'_1\,\alpha'_2\,\alpha'_3\,\alpha'_4\,a,\\
 p(b,d)=b\,\alpha'_3\,\alpha'_4\,\beta_1\,\beta_2\,d, &
 q(b,d)=b\,\alpha_3\,\alpha_4\,\beta'_1\,\beta'_2\,d,\\
 p(e,f)=e\,\beta'_1\,\beta'_2\,\alpha'_1\,\alpha'_2\,f, &
 q(g,h)=g\,\beta_1\,\beta_2\,\alpha_1\,\alpha_2\,h.
\end{array}
\]
Each has six vertices. The three \(p\)-paths are pairwise vertex-disjoint,
the three \(q\)-paths are pairwise vertex-disjoint, and \(\mu\) exchanges
\(p(c,a)\) with \(q(c,a)\), \(p(b,d)\) with \(q(b,d)\), and \(p(e,f)\) with
\(q(g,h)\).
\end{hbdefinition}

The figure labels edges rather than vertices: \(c\to\alpha_1\),
\(\alpha_1\to\alpha_2\), \(\alpha_2\to\alpha_3\), \(\alpha_3\to\alpha_4\),
\(\alpha_4\to a\), \(b\to\alpha_3\), \(\alpha_4\to\beta'_1\),
\(g\to\beta_1\), \(\beta_1\to\beta_2\), \(\beta_2\to\alpha_1\),
\(\alpha_2\to h\), \(\beta_2\to d\) carry the labels \(5,4,3,2,1,6,7,8,9,
10,11,12\), and the primed edges the primed labels, while the four edges that
join a switch to the rest of a graph are labelled \(A,B,C,D\) at the
vertices called \(a,b,c,d\) here
\cite[Fig.~9.1, p.~479]{BangJensenGutin2007DigraphsCorrected}. Kolaitis and
Vardi list the six paths as sequences of these labels; their lists for
\(p(b,d)\) and \(q(b,d)\) read \(6',2',7,9,12\) and \(6,2,7',9',12'\)
\cite[p.~30]{KolaitisVardi1995DatalogManuscript}, whereas in the figure the
edge leaving the head of \(2'\) is \(7'\) and the edge leaving the head of
\(2\) is \(7\). The vertex sequences above follow the figure; with the labels
\(7\) and \(7'\) as printed, the two lists would not be paths.

\begin{hblemma}{thm:p03-w03-p03-061-switch}{proved=annex:p03-w03-p03-061-switch}{Paths through a switch}
Let \(P\) and \(Q\) be paths passing through \(\mathbf S\) with no vertex in
common, where \(P\) ends at \(a\) and \(Q\) starts at \(b\). Then either
\(P=p(c,a)\) and \(Q=p(b,d)\), and \(p(e,f)\) is the only path passing
through \(\mathbf S\) that has no vertex in common with \(P\) or \(Q\); or
\(P=q(c,a)\) and \(Q=q(b,d)\), and \(q(g,h)\) is the only such path.
\end{hblemma}

This is Kolaitis and Vardi's Lemma~6.4, whose proof they leave to an
inspection of the switch \cite[Lemma~6.4, p.~30]{KolaitisVardi1995DatalogManuscript};
Bang-Jensen and Gutin state it as their Lemma~9.2.4 and leave its proof as an
exercise \cite[Lemma~9.2.4, p.~478]{BangJensenGutin2007DigraphsCorrected}.
In the graphs below the first alternative will mean that the literal of the
switch is true, the second that it is false.

\begin{hbdefinition}{def:p03-w03-p03-061-fhw-graph}{The graph \(\mathbf G_\varphi\) of a formula and the complete formulas \(\varphi_k\)}
Let \(\varphi\) be a formula in conjunctive normal form with variables
\(x_1,\ldots,x_k\) (\(k\ge1\)) and clauses \(c_1,\ldots,c_l\) (\(l\ge1\)),
each clause a nonempty set of literals; \(\bar x_i\) is the negation of
\(x_i\). An \emph{occurrence} is a pair \(o=(y,j)\) with \(y\in c_j\); its
literal is \(y\) and its variable is the variable of \(y\). Fix an
enumeration \(o_1,\ldots,o_N\) of the occurrences (\(N\ge1\)), and for each
literal \(y\) let \(o^y_1,\ldots,o^y_{m(y)}\) be the occurrences of \(y\)
in increasing order of the clause index. The directed graph
\(\mathbf G_\varphi\) has the vertices \(z_1,z_2,z_3,z_4\),
\(u_1,\ldots,u_k\), \(v_1,\ldots,v_k\), \(n_0,\ldots,n_l\), and for each
occurrence \(o\) and each vertex \(w\) of \(\mathbf S\) a vertex
\(w\langle o\rangle\) (the \emph{switch of} \(o\)). Its edges are:
\begin{enumerate}
\item \(w\langle o\rangle\to w'\langle o\rangle\) for every occurrence
      \(o\) and every edge \(w\to w'\) of \(\mathbf S\);
\item \(d\langle o_r\rangle\to b\langle o_{r+1}\rangle\) and
      \(a\langle o_{r+1}\rangle\to c\langle o_r\rangle\) for
      \(1\le r<N\);
\item \(z_1\to c\langle o_N\rangle\), \(a\langle o_1\rangle\to z_2\),
      \(z_3\to b\langle o_1\rangle\), \(d\langle o_N\rangle\to u_1\) and
      \(n_l\to z_4\);
\item for each \(i\) and each literal \(y\in\{x_i,\bar x_i\}\): if
      \(m(y)\ge1\), the edges \(u_i\to g\langle o^y_1\rangle\),
      \(h\langle o^y_s\rangle\to g\langle o^y_{s+1}\rangle\) for
      \(1\le s<m(y)\), and \(h\langle o^y_{m(y)}\rangle\to v_i\) (the
      \emph{column} of \(y\)); if \(m(y)=0\), the edge \(u_i\to v_i\);
\item \(v_i\to u_{i+1}\) for \(1\le i<k\), and \(v_k\to n_0\);
\item \(n_{j-1}\to e\langle o\rangle\) and \(f\langle o\rangle\to n_j\) for
      every occurrence \(o=(y,j)\).
\end{enumerate}
\(\mathbf B_\varphi\) is the \(\sigma_{\mathbf H_1}\)-structure with graph
\(\mathbf G_\varphi\) and \(d_{s_1},d_{s_2},d_{s_3},d_{s_4}\) interpreted by
\(z_1,z_2,z_3,z_4\); it belongs to \(\mathcal F_{\mathbf H_1}\).

For \(k\ge1\) the \emph{complete formula} \(\varphi_k\) on
\(x_1,\ldots,x_k\) has as clauses the \(2^k\) sets that contain exactly one of
\(x_i,\bar x_i\) for each \(i\) \cite[p.~34]{KolaitisVardi1995DatalogManuscript}.
Every literal occurs in \(2^{k-1}\) clauses, and \(\varphi_k\) is
unsatisfiable: under an assignment, the clause made of the literals it makes
false is false.
\end{hbdefinition}

The four groups of edges in items~2--6 are Kolaitis and Vardi's steps: each
occurrence has its own switch, whose \(q(g,h)\) replaces an edge of the column
of its literal and whose \(p(e,f)\) is one of the paths from \(n_{j-1}\) to
\(n_j\); the switches are chained \(d\) to \(b\) and \(a\) to \(c\); the
bottom of the last variable block leads to \(n_0\); and four terminal
vertices are attached \cite[pp.~30--31]{KolaitisVardi1995DatalogManuscript}.
They call the terminals \(s_1,\ldots,s_4\), here the names of the nodes of
\(\mathbf H_1\). Their display of \(\varphi_2\) lists the clause
\(\bar x_1\vee\bar x_2\) twice and omits \(x_1\vee\bar x_2\)
\cite[p.~34]{KolaitisVardi1995DatalogManuscript}; the definition above is the
one their description states. Kolaitis's hint names the reduction from
3-satisfiability \cite[p.~79]{Kolaitis2003ExpressivePowerNotes}; the
argument needs clauses of length \(k\), so it uses the reduction for
arbitrary formulas in conjunctive normal form, as Kolaitis and Vardi do.

\begin{hbdefinition}{def:p03-w03-p03-061-templates}{Modes, templates and good sets}
Let \(\varphi\) and \(\mathbf G_\varphi\) be as in
\cref{def:p03-w03-p03-061-fhw-graph}. An \emph{assignment} \(\nu\) maps each
\(x_i\) to \(0\) or \(1\); put \(\nu(\bar x_i)\defeq1-\nu(x_i)\), and let
\(\lambda^\nu_i\) be the literal among \(x_i,\bar x_i\) that \(\nu\) makes
false. The \emph{mode} of an occurrence \(o\) under \(\nu\) is \(p\) if
\(\nu\) makes its literal true and \(q\) otherwise. A \emph{clause choice}
\(\kappa\) picks for each \(j\) an occurrence \(\kappa(j)\) of the form
\((y,j)\). Write \(\mathrm{CA}_\nu(o)\) for the vertex sequence
\(p(c,a)\langle o\rangle\) or \(q(c,a)\langle o\rangle\) according to the
mode of \(o\), and \(\mathrm{BD}_\nu(o)\) likewise for \(p(b,d)\) and
\(q(b,d)\). The \emph{templates} are the concatenated vertex sequences
\[
\begin{aligned}
 T_1(\nu)&\defeq z_1,\ \mathrm{CA}_\nu(o_N),\ \mathrm{CA}_\nu(o_{N-1}),\ \ldots,\\
 &\qquad\mathrm{CA}_\nu(o_1),\ z_2,\\
 T_2(\nu,\kappa)&\defeq z_3,\ \mathrm{BD}_\nu(o_1),\ \ldots,\ \mathrm{BD}_\nu(o_N),\\
 &\qquad\Gamma_1,\ \ldots,\ \Gamma_k,\ n_0,\\
 &\qquad p(e,f)\langle\kappa(1)\rangle,\ n_1,\ \ldots,\\
 &\qquad p(e,f)\langle\kappa(l)\rangle,\ n_l,\ z_4,
\end{aligned}
\]
where, with \(y=\lambda^\nu_i\), \(\Gamma_i\) is \(u_i\), then
\(q(g,h)\langle o^y_1\rangle,\ldots,q(g,h)\langle o^y_{m(y)}\rangle\), then
\(v_i\). The formula is \emph{balanced} when \(m(x_i)=m(\bar x_i)\) for every
\(i\); then the lengths
\[
 L_1\defeq6N+2,\qquad
 L_2\defeq6N+\textstyle\sum_{i=1}^k\bigl(6m(x_i)+2\bigr)+7l+3
\]
of \(T_1(\nu)\) and \(T_2(\nu,\kappa)\) do not depend on \(\nu,\kappa\). The
positions of a template are numbered from \(0\); \(T[t]\) is the entry at
position \(t\). The \emph{segment of clause} \(j\) is the set of the six
positions of \(T_2\) occupied by \(p(e,f)\langle\kappa(j)\rangle\). A set
\(D\) of positions of \(T_2\) is \emph{good} for \((\nu,\kappa)\) when
\(\nu\) makes the literal of \(\kappa(j)\) true for every \(j\) whose segment
meets \(D\).
\end{hbdefinition}

\begin{hblemma}{thm:p03-w03-p03-061-templates}{proved=annex:p03-w03-p03-061-templates}{Templates are walks that are injective on good positions}
Let \(\varphi\) be a formula as in \cref{def:p03-w03-p03-061-fhw-graph},
\(\nu\) an assignment and \(\kappa\) a clause choice.
\begin{enumerate}
\item Any two consecutive entries of \(T_1(\nu)\) or of \(T_2(\nu,\kappa)\)
      are joined by an edge of \(\mathbf G_\varphi\) directed from the
      earlier to the later entry.
\item If \(D\) is a set of positions of \(T_2\) that is good for
      \((\nu,\kappa)\), the entries of \(T_1(\nu)\) at all its positions and
      the entries of \(T_2(\nu,\kappa)\) at the positions outside every
      clause segment and at the positions in \(D\) are pairwise distinct
      vertices (no vertex occurs at two of these positions).
\item Let \(\varphi\) be balanced. The entry of \(T_1(\nu)\) at a position
      inside \(\mathrm{CA}_\nu(o_r)\) other than its first and last, and the
      entry of \(T_2(\nu,\kappa)\) at such a position inside
      \(\mathrm{BD}_\nu(o_r)\), depends only on the value of \(\nu\) at the
      variable of \(o_r\); the entry of \(T_2\) at a position strictly between
      those of \(u_i\) and \(v_i\) depends only on \(\nu(x_i)\); the entry
      at a position in the segment of clause \(j\) depends only on
      \(\kappa(j)\); every other entry is the same for all \(\nu,\kappa\).
\end{enumerate}
\end{hblemma}

\begin{hbtheorem}{thm:p03-w03-p03-061-fhw}{proved=annex:p03-w03-p03-061-fhw}{The reduction of Fortune, Hopcroft and Wyllie}
Let \(\varphi\) be a formula as in \cref{def:p03-w03-p03-061-fhw-graph}.
\(\mathbf G_\varphi\) has simple paths from \(z_1\) to \(z_2\) and from
\(z_3\) to \(z_4\) with no vertex in common, that is
\(\mathrm{FiSH}(\mathbf H_1)(\mathbf B_\varphi)=1\), if and only if
\(\varphi\) is satisfiable. In particular
\(\mathrm{FiSH}(\mathbf H_1)(\mathbf B_{\varphi_k})=0\) for every \(k\ge1\).
\end{hbtheorem}

Kolaitis and Vardi attribute the equivalence to Fortune, Hopcroft and Wyllie
and prove both directions by tracing the two paths through the switches
\cite[pp.~31--33]{KolaitisVardi1995DatalogManuscript}. Only the direction
from paths to a satisfying assignment is used for the lower bound; the
annex proves both, the converse from \cref{thm:p03-w03-p03-061-templates}.

\begin{hbdefinition}{def:p03-w03-p03-061-ab}{The graphs \(\mathbf A_k\) and \(\mathbf B_k\)}
Let \(k\ge1\), let \(\mathbf B_k\defeq\mathbf B_{\varphi_k}\), and let
\(L_1,L_2\) be the template lengths of the balanced formula \(\varphi_k\).
\(\mathbf A_k\) is the \(\sigma_{\mathbf H_1}\)-structure whose graph is the
disjoint union of the paths \(\eta_0\to\eta_1\to\cdots\to\eta_{L_1-1}\) and
\(\eta'_0\to\eta'_1\to\cdots\to\eta'_{L_2-1}\), with \(d_{s_1},d_{s_2},d_{s_3},
d_{s_4}\) interpreted by \(\eta_0,\eta_{L_1-1},\eta'_0,\eta'_{L_2-1}\)
\cite[pp.~34--35]{KolaitisVardi1995DatalogManuscript}. Let \(W\) be the set
of these four vertices. For an assignment \(\nu\) and a clause choice
\(\kappa\), \(h_{\nu,\kappa}\) maps \(\eta_t\) to \(T_1(\nu)[t]\) and
\(\eta'_t\) to \(T_2(\nu,\kappa)[t]\). A set \(D\) of vertices of
\(\mathbf A_k\) is good for \((\nu,\kappa)\) when
\(\{t:\eta'_t\in D\}\) is. Finally \(\mathcal I_k\) is the set of all maps
\(h_{\nu,\kappa}\restriction(D\cup W)\) with \(D\subseteq A_k-W\),
\(|D|\le k\) and \(D\) good for \((\nu,\kappa)\).
\end{hbdefinition}

\begin{hbtheorem}{thm:p03-w03-p03-061-family}{proved=annex:p03-w03-p03-061-family}{Duplicator's winning family}
For every \(k\ge1\),
\[
 \mathrm{FiSH}(\mathbf H_1)(\mathbf A_k)=1,\qquad
 \mathrm{FiSH}(\mathbf H_1)(\mathbf B_k)=0,
\]
and \(\mathcal I_k\) is a winning family for the \((\exists,\neq,k)\)-pebble game on \(\mathbf A_k\)
and \(\mathbf B_k\).
\end{hbtheorem}

A member of \(\mathcal I_k\) records Duplicator's answers to the pebbles in
\(D\). A pebble on the first path is answered inside \(p(c,a)\) or
\(q(c,a)\) of a switch, as the truth value of its literal dictates; a pebble
in the first part of the second path inside \(p(b,d)\) or \(q(b,d)\), by the
same rule; a pebble in the block of \(x_i\) in the column of the false
literal; and a pebble in the segment of clause \(j\) inside \(p(e,f)\) of an
occurrence whose literal she makes true. The forth step needs a literal of
the clause whose variable no pebble has fixed yet, which exists because at
most \(k-1\) other pebbles fix at most \(k-1\) variables and every clause of
\(\varphi_k\) mentions all \(k\) variables. These are Kolaitis and Vardi's
four cases \cite[pp.~36--38]{KolaitisVardi1995DatalogManuscript}.

\begin{hbtheorem}{thm:p03-w03-p03-061-h1}{proved=annex:p03-w03-p03-061-h1}{The 2-disjoint paths query is not existentially definable}
For no \(k\ge1\) is \(\mathrm{FiSH}(\mathbf H_1)\)
\(\exists\mathrm L^k_{\infty\omega}(\neq)\)-definable on
\(\mathcal F_{\mathbf H_1}\). Hence \(\mathrm{FiSH}(\mathbf H_1)\) is not
\(\exists\mathrm L^\omega_{\infty\omega}(\neq)\)-definable, and no
Datalog\((\neq)\) program defines it.
\end{hbtheorem}

This is Kolaitis and Vardi's Theorem~6.6
\cite[Thm~6.6, p.~34]{KolaitisVardi1995DatalogManuscript}.

\begin{hblemma}{thm:p03-w03-p03-061-quotient}{proved=annex:p03-w03-p03-061-quotient}{Identifying distinguished vertices}
Let \(\mathbf A,\mathbf B\) be directed graphs with distinguished vertices
\(a_1,\ldots,a_n\) of \(\mathbf A\) and \(b_1,\ldots,b_n\) of \(\mathbf B\),
each list without repetitions, read as structures with constants
\(c_1,\ldots,c_n\). Let \(\equiv\) be an equivalence relation on
\(\{1,\ldots,n\}\). Let \(\mathbf A/{\equiv}\) arise from \(\mathbf A\) by
identifying \(a_i\) with \(a_j\) whenever \(i\equiv j\): its vertices are
the classes \([x]\) (singletons except for the identified distinguished
vertices), its edges are the pairs \(([x],[x'])\) with \(x\to x'\) in
\(\mathbf A\), and \(c_i\) denotes \([a_i]\). Define \(\mathbf B/{\equiv}\)
in the same way. If \(\mathcal I\) is a winning family for the
\((\exists,\neq,k)\)-pebble game on \(\mathbf A\) and \(\mathbf B\), then
\(\{f/{\equiv}:f\in\mathcal I\}\), where
\(f/{\equiv}\defeq\{([x],[y]):(x,y)\in f\}\), is a winning family for that
game on \(\mathbf A/{\equiv}\) and \(\mathbf B/{\equiv}\).
\end{hblemma}

\begin{hbtheorem}{thm:p03-w03-p03-061-h2h3}{proved=annex:p03-w03-p03-061-h2h3}{The basic patterns \(\mathbf H_2\) and \(\mathbf H_3\)}
For no \(k\ge1\) is \(\mathrm{FiSH}(\mathbf H_2)\)
\(\exists\mathrm L^k_{\infty\omega}(\neq)\)-definable on
\(\mathcal F_{\mathbf H_2}\), and for no \(k\ge1\) is
\(\mathrm{FiSH}(\mathbf H_3)\) \(\exists\mathrm L^k_{\infty\omega}(\neq)\)-definable
on \(\mathcal F_{\mathbf H_3}\). The witnesses are obtained from
\(\mathbf A_k,\mathbf B_k\) by identifying the second and third
distinguished vertices (for \(\mathbf H_2\)), respectively the first with the
fourth and the second with the third (for \(\mathbf H_3\)).
\end{hbtheorem}

These are the modifications Kolaitis and Vardi indicate in the proof of their
Theorem~6.7 \cite[Thm~6.7, pp.~38--39]{KolaitisVardi1995DatalogManuscript},
where they then play the game on the modified graphs; the transfer of the
winning family is \cref{thm:p03-w03-p03-061-quotient}.

\begin{hbtheorem}{thm:p03-w03-p03-061-dichotomy}{proved=annex:p03-w03-p03-061-dichotomy}{Descriptive dichotomy for fixed subgraph homeomorphism queries}
Let \(\mathbf H\) be a pattern graph in the sense of
\cref{def:p03-w03-p03-061-pattern}.
\begin{enumerate}
\item If \(\mathbf H\) is a star graph with \(m\) edges,
      \(\mathrm{FiSH}(\mathbf H)\) is
      \(\exists\mathrm L^{m+1}_{\infty\omega}(\neq)\)-definable on
      \(\mathcal F_{\mathbf H}\) and is defined there by a Datalog\((\neq)\)
      program.
\item If \(\mathbf H\) is not a star graph, \(\mathrm{FiSH}(\mathbf H)\) is
      not \(\exists\mathrm L^k_{\infty\omega}(\neq)\)-definable on
      \(\mathcal F_{\mathbf H}\) for any \(k\), hence not
      \(\exists\mathrm L^\omega_{\infty\omega}(\neq)\)-definable and not
      definable by any Datalog\((\neq)\) program.
\end{enumerate}
\end{hbtheorem}

This is Kolaitis's Theorem~8.19 for loopless patterns without isolated
nodes \cite[Thm~8.19, p.~79]{Kolaitis2003ExpressivePowerNotes}; for
isolated nodes see \cref{def:p03-w03-p03-061-edgeless}.

\begin{hbexample}{def:p03-w03-p03-061-phi1}{The case \(k=1\)}
The formula \(\varphi_1\) has the clauses \(c_1=\{x_1\}\) and
\(c_2=\{\bar x_1\}\), hence the occurrences \(o_1=(x_1,1)\) and
\(o_2=(\bar x_1,2)\), \(N=2\) and \(m(x_1)=m(\bar x_1)=1\). In
\(\mathbf G_{\varphi_1}\) the column of \(x_1\) is
\(u_1\to q(g,h)\langle o_1\rangle\to v_1\) and the column of \(\bar x_1\) is
\(u_1\to q(g,h)\langle o_2\rangle\to v_1\); the paths from \(n_0\) to \(n_1\)
and from \(n_1\) to \(n_2\) run through \(p(e,f)\langle o_1\rangle\) and
\(p(e,f)\langle o_2\rangle\). Two disjoint paths from \(z_1\) to \(z_2\) and
from \(z_3\) to \(z_4\) would have to fix a mode in each switch, use one
column, which forces the switch of that column into mode \(q\), and pass both
clause paths, which forces both switches into mode \(p\); so there are none.
Here \(L_1=14\) and \(L_2=37\). In the \((\exists,\neq,1)\)-pebble game
Spoiler's single pebble on \(\mathbf A_1\) is answered by
\(h_{\nu,\kappa}\) for a pair \((\nu,\kappa)\) that makes its position good:
a pebble in the segment of clause \(1\) is answered with \(\nu(x_1)=1\), one in
the segment of clause \(2\) with \(\nu(x_1)=0\). Two pebbles defeat this
answer scheme: once a pebble strictly between \(\eta'_{13}\) and
\(\eta'_{20}\), which are sent to \(u_1\) and \(v_1\), has been answered in the
column of \(x_1\), so that \(\nu(x_1)=0\), a second pebble in the segment of
clause \(1\) has no answer \(h_{\nu',\kappa'}\) agreeing with the first that
makes both positions good, because clause \(1\) contains only \(x_1\). This
is where the bound on the number of pebbles enters the family.
\end{hbexample}

\begin{hbwarning}{def:p03-w03-p03-061-family-limit}{What the construction does and does not give}
The family \(\mathcal I_k\) shows that Duplicator wins with \(k\) pebbles on
\(\mathbf A_k\) and \(\mathbf B_k\); it says nothing about \(k+1\) pebbles on
the same pair, and the example above shows that \(\mathcal I_k\) itself
cannot be enlarged to positions with \(k+1\) pebbles. The argument proves
non-definability in the existential logic with inequalities, and through
\cref{thm:p03-w03-p03-061-datalog} in Datalog\((\neq)\); it gives no
information about \(\mathrm L^\omega_{\infty\omega}\), where Spoiler may also
play in \(\mathbf B_k\). Kolaitis records that case as open in the
2003 manuscript \cite[p.~79]{Kolaitis2003ExpressivePowerNotes}; that
historical report is not a claim about its subsequent status. The graphs \(\mathbf B_k\)
have more than \(20k2^k\) vertices; the construction concerns definability,
not the cost of evaluating the query.
\end{hbwarning}

\section{Relationships, limits and sources}
\label{sec:p03-w03-p03-061-relations}

The existential game refines the \(k\)-pebble game of
\cref{def:p03-w03-p03-055-pebble}, which is set up there for relational
vocabularies.

\begin{hblemma}{thm:p03-w03-p03-061-refines}{proved=annex:p03-w03-p03-061-refines}{A win in the \(k\)-pebble game is a win in the existential games}
Let \(\sigma\) be a relational vocabulary, \(k\ge1\), and \(\mathbf A,\mathbf
B\) \(\sigma\)-structures. Every winning family for the \(k\)-pebble game of
\cref{def:p03-w03-p03-055-pebble} on \(\mathbf A\) and \(\mathbf B\) is a
winning family for the \((\exists,\neq,k)\)-pebble game and for the
\((\exists,k)\)-pebble game on \(\mathbf A\) and \(\mathbf B\). The converse
fails: by \cref{thm:p03-w03-p03-061-cliques}, Duplicator wins the
\((\exists,\neq,j)\)-pebble games on \(\mathbf K_k\) and \(\mathbf K_{k+1}\)
for all \(j\), while Spoiler wins the \((k+1)\)-pebble game on them.
\end{hblemma}

Kolaitis motivates the existential game by this restriction of Spoiler's
moves \cite[p.~71]{Kolaitis2003ExpressivePowerNotes}. The criterion
\cref{thm:p03-w03-p03-061-criterion} is the existential counterpart of the
definability criterion for \(\mathrm L^k_{\infty\omega}\) proved in
\cref{sec:p03-w03-p03-055-pebble}. By \cref{thm:p03-w03-p03-061-homomorphism},
a homomorphism from \(\mathbf A\) to \(\mathbf B\) is the strongest form of
Duplicator's win; for structures given as canonical databases of conjunctive
queries this is the Chandra--Merlin test (\cref{thm:p09-chandra-merlin}).

\emph{Limits.} The criterion characterizes
\(\exists\mathrm L^\omega_{\infty\omega}\), which is strictly larger than
Datalog on finite structures (\cref{thm:p03-w03-p03-061-proper}); a query can be
\(\exists\mathrm L^\omega_{\infty\omega}\)-definable and still not
expressible in Datalog. Kolaitis records as open whether the non-star
\(\mathrm{FiSH}(\mathbf H)\) queries fail to be definable even in
\(\mathrm L^\omega_{\infty\omega}\), the critical case being the 2-disjoint
paths query \cite[p.~79]{Kolaitis2003ExpressivePowerNotes}. This is the
status reported in that 2003 source; the arguments here establish no
stronger non-definability theorem. On undirected graphs every
\(\mathrm{FiSH}(\mathbf H)\) is polynomial-time solvable by the work of
Robertson and Seymour, and Barland showed \(\mathrm{FiSH}(\mathbf C_3)\)
definable in least fixed-point logic on finite undirected graphs, as Kolaitis
reports; any descriptive dichotomy there would have a different boundary
\cite[p.~79]{Kolaitis2003ExpressivePowerNotes}.

\emph{Sources.} Definitions and statements: Kolaitis's manuscript of
1 August 2003, Section~8 \cite[pp.~69--79]{Kolaitis2003ExpressivePowerNotes}.
Conventions on pattern graphs, the program for star patterns, the
reduction to basic patterns, the reduction of Fortune, Hopcroft and Wyllie
and the graphs \(\mathbf A_k,\mathbf B_k\): the Kolaitis--Vardi manuscript
\cite[pp.~22--39]{KolaitisVardi1995DatalogManuscript}. The switch: the figure
in Bang-Jensen and Gutin \cite[Fig.~9.1, p.~479]{BangJensenGutin2007DigraphsCorrected}.
The proofs in the annex are written for this chapter.
